% lecture_8.tex

\section*{Cursul VIII}

\begin{exercitiu}{1}
    Sa se reia analiza aproximarii cu diferente finite centrale a lui $ J_F(x\p{k})$ cu $ J_{\tilde{h}}\p{k}$ referitoare la:
    \begin{itemize}
        \item[-] corectitudinea aproximarii $J_{\tilde{h}}\p{k}$, 
        \item[-] normalizarea lui $J_{\tilde{h}}\p{k}$,
        \item[-] analiza erorii de aproximare.
    \end{itemize}
\end{exercitiu}

\begin{demonstratie}

    \noindent- corectitudinea
    \begin{align*}
        J_F(x\p{k}) w &= J_F(x\p{k}) \left( \sum_{j=1}^n w_j e\p{j} \right) \\
        &= \sum_{j=1}^n w_j \left[ J_F(x\p{k}) e\p{j} \right] \\
        &= \sum_{j=1}^n w_j \frac{\partial F}{\partial x_j}(x\p{k}) \\
        &\simeq \sum_{j=1}^n w_j \frac{1}{2\tilde{h}_j\p{k}} \left[ F(x\p{k} + \tilde{h}_j\p{k} e\p{j}) - F(x\p{k} - \tilde{h}_j\p{k} e\p{j}) \right] \\
        &= \sum_{j=1}^n w_j \left[ \tilde{J}_{\tilde{h}}\p{k} e\p{j} \right] \\
        &= \tilde{J}_{\tilde{h}}\p{k} \left( \sum_{j=1}^n w_j e\p{j} \right) \\
        &= J_{\tilde{h}}\p{k} w
    \end{align*}

    In consecinta, $ J_F(x\p{k}) w \simeq J_{\tilde{h}}\p{k} w, \quad \forall w \in \R^n$.

    \bigskip

    \noindent- normalizarea
    $$
    \frac{\partial F}{\partial x_j}(x\p{k}) \simeq
    \begin{cases}
        \displaystyle \frac{1}{2\tilde{h}_j\p{k} \|x\p{k}\|} \left[ F(x\p{k} + \tilde{h}_j\p{k} e\p{j}) - F(x\p{k} - \tilde{h}_j\p{k} e\p{j}) \right], & \text{dacă } x\p{k} \neq 0 \\[15pt]
        \displaystyle \frac{1}{2\tilde{h}_j\p{k}} \left[ F(x\p{k} + \tilde{h}_j\p{k} e\p{j}) - F(x\p{k} - \tilde{h}_j\p{k} e\p{j}) \right], & \text{dacă } x\p{k} = 0
    \end{cases}
    $$

    \newpage

    \noindent- analiza erorii

    \begin{align*}
        (\tilde{J}_{\tilde{h}}\p{k})_{\cdot, j} &= \frac{\tilde{F}(x\p{k} + \tilde{h}_j\p{k} \|x\p{k}\| e\p{j}) - \tilde{F}(x\p{k} - \tilde{h}_j\p{k} \|x\p{k}\| e\p{j})}{2\tilde{h}_j\p{k} \|x\p{k}\|} \\
        &= \frac{F(x\p{k} + \tilde{h}_j\p{k} \|x\p{k}\| e\p{j}) - F(x\p{k} - \tilde{h}_j\p{k} \|x\p{k}\| e\p{j})}{2\tilde{h}_j\p{k} \|x\p{k}\|} \\
        &\quad + \frac{\mathcal{E}(x\p{k} + \tilde{h}_j\p{k} \|x\p{k}\| e\p{j}) - \mathcal{E}(x\p{k} - \tilde{h}_j\p{k} \|x\p{k}\| e\p{j})}{2\tilde{h}_j\p{k} \|x\p{k}\|}
    \end{align*}

    Eroarea totala a aproximarii este:

    \begin{align*}
        \| J_F(x\p{k})e\p{j} - (\tilde{J}_{\tilde{h}}\p{k})_{\cdot, j} \| &\leq \left\| \frac{\partial F}{\partial x_j}(x\p{k}) - \frac{F(x\p{k} + \dots) - F(x\p{k} - \dots)}{2\tilde{h}_j\p{k} \|x\p{k}\|} \right\| \\
        &\quad + \left\| \frac{\mathcal{E}(x\p{k} + \dots) - \mathcal{E}(x\p{k} - \dots)}{2\tilde{h}_j\p{k} \|x\p{k}\|} \right\| \\
        &= O\left( (\tilde{h}_j\p{k})^2 \|x\p{k}\|^2 \right) + O\left( \frac{\epsilon_M}{\tilde{h}_j\p{k} \|x\p{k}\|} \right).
    \end{align*} 

\end{demonstratie}

\begin{lema}{Sherman-Morrison}
    Fie $A \in \M_n(\R)$ inversabila si $u, v \in \R^n$. Atunci:

\begin{itemize}
    \item[(i)] $A + u v^T \in \M_n(\R)$ inversabila $\iff 1 + v^T A^{-1} u \neq 0$,
    \item[(ii)] $1 + v^T A^{-1} u \neq 0 \implies$
    $$ (A + u v^T)^{-1} = A^{-1} - \frac{A^{-1} u v^T A^{-1}}{1 + v^T A^{-1} u}. $$
\end{itemize}
\end{lema}

\begin{demonstratie}

    \noindent i)

    Se considera $ w \in \R^{n}$ a.i. $ w = A^{-1} u$. Atunci:

    $$
    A + u v^{T} = A (I_n + A^{-1} u v^{T}) = A(I_n + w v^{T}).
    $$

    Intrucat produsul dintre doua matrici inversabile este o matrice inversabila, demonstratia se reduce la inversabilitatea matricii $ B = I_n + w v^{T}$. Pentru $ x \in \R^{n}$ avem:
    $$
    B x = (I_n + w v^{T}) x = x + w v^{T} x.
    $$

    Cazul $ x = w $:
    $$
    B w = w + w v^{T} w = w(1 + v^{T} w),
    $$
    ceea ce inseamna ca $ w $ este vectorul propriu al matricii $ B $ cu valoarea proprie asociata $\lambda_w = (1 + v^{T} w)$.

    Cazul $ x = z : v^{T} z = 0 $:
    $$
    B z = z + w v^{T} z = z,
    $$
    ceea ce inseamna ca pe directiile ortogonale lui $ v $, directia e invarianta la transformarea $ B \implies \lambda_i = 1 $, $ \forall i \in \{1, \dots, n - 1\}$.

    In consecinta, ca toate valorile proprii ale matricii B sa fie nenule, trebuie ca $ \lambda_w = (1 + v^{T} w) \neq 0 $. Prin urmare, inversabilitatea matricii $ B $, si implicit, a matricii $ A $ depinde de conditia ca $1 + v^T A^{-1} u \neq 0$. 

    \bigskip

    \noindent ii)

    Pentru $ X , Y \in \M_n(\R) $, daca $ X Y = I_n \implies Y = X^{-1} $. Atunci: 
    \begin{align*}
        (A + uv^T) & \left( A^{-1} - \frac{A^{-1}uv^TA^{-1}}{1 + v^TA^{-1}u} \right) \\
        &= AA^{-1} - A \frac{A^{-1}uv^TA^{-1}}{1 + v^TA^{-1}u} + uv^TA^{-1} - uv^T \frac{A^{-1}uv^TA^{-1}}{1 + v^TA^{-1}u} \\
        &= I_n - \frac{uv^TA^{-1}}{1 + v^TA^{-1}u} + uv^TA^{-1} - \frac{u(v^TA^{-1}u)v^TA^{-1}}{1 + v^TA^{-1}u} \\
        &= I_n + uv^TA^{-1} - \frac{uv^TA^{-1} + u(v^TA^{-1}u)v^TA^{-1}}{1 + v^TA^{-1}u} \\
        &= I_n + uv^TA^{-1} - \frac{u \left( 1 + v^TA^{-1}u \right) v^TA^{-1}}{1 + v^TA^{-1}u} \\
        &= I_n + uv^TA^{-1} - uv^TA^{-1} \\
        &= I_n
    \end{align*}
\end{demonstratie}